\documentclass{article}

% if you need to pass options to natbib, use, e.g.:
%     \PassOptionsToPackage{numbers, compress}{natbib}
% before loading neurips_2026

% The authors should use one of these tracks.
% Before accepting by the NeurIPS conference, select one of the options below.
% 0. "default" for submission
\usepackage[preprint]{neurips_2026}

\usepackage[utf8]{inputenc} % allow utf-8 input
\usepackage[T1]{fontenc}    % use 8-bit T1 fonts
\usepackage{hyperref}       % hyperlinks
\usepackage{url}            % simple URL typesetting
\usepackage{booktabs}       % professional-quality tables
\usepackage{amsfonts}       % blackboard math symbols
\usepackage{nicefrac}       % compact symbols for 1/2, etc.
\usepackage{microtype}      % microtypography
\usepackage{xcolor}         % colors
\usepackage{graphicx}
\usepackage{booktabs}
\usepackage{multirow}

% Note. For the workshop paper template, both \title{} and \workshoptitle{} are required, with the former indicating the paper title shown in the title and the latter indicating the workshop title displayed in the footnote. 
\title{Exploring Deep Learning Training Paradigms and Model Architectures}


% The \author macro works with any number of authors. There are two commands
% used to separate the names and addresses of multiple authors: \And and \AND.
%
% Using \And between authors leaves it to LaTeX to determine where to break the
% lines. Using \AND forces a line break at that point. So, if LaTeX puts 3 of 4
% authors names on the first line, and the last on the second line, try using
% \AND instead of \And before the third author name.


\author{%
    M. Ibrahim\\
    \texttt{https://github.com/Ibrahim778/ATML\_PA1}
}


\begin{document}


\maketitle


\begin{abstract}
    This project explores the what and why behind the architectural decisions and training methods of several different deep learning models.
\end{abstract}



\section{Convolutional Neural Networks (CNN) Training Comparisons}

% requires \usepackage{booktabs}

% ---------- Table 1: Clean baseline ----------
\begin{table}[h]
\centering
\caption{Task 1: Clean baseline performance on the held-out test subset ($n=240$).}
\label{tab:t1-clean-baseline}
\begin{tabular}{lccc}
\toprule
Model & Accuracy & Macro-F1 & Mean max.\ confidence \\
\midrule
ResNet-50          & 0.972 & 0.972 & 0.947 \\
ViT-B/16            & 0.970 & 0.970 & 0.965 \\
CLIP ViT-B/32 (linear probe) & 0.966 & 0.966 & 0.241 \\
CLIP ViT-B/32 (zero-shot)    & 0.936 & 0.936 & 0.925 \\
\bottomrule
\end{tabular}
\end{table}


% ---------- Table 2: Color-bias interventions ----------
\begin{table}[h]
\centering
\caption{Task 1: Effect of color interventions (grayscale, $90^\circ$ hue rotation) on accuracy, prediction consistency vs.\ the clean prediction, and mean max.\ confidence.}
\label{tab:t1-color-bias}
\begin{tabular}{llccc}
\toprule
Model & Intervention & Accuracy & Consistency vs.\ clean & Mean max.\ confidence \\
\midrule
\multirow{2}{*}{ResNet-50}      & Grayscale      & 0.946 & 0.952 & 0.943 \\
                                 & Hue rot.\ $90^\circ$ & 0.938 & 0.954 & 0.927 \\
\addlinespace
\multirow{2}{*}{ViT-B/16}        & Grayscale      & 0.950 & 0.978 & 0.955 \\
                                 & Hue rot.\ $90^\circ$ & 0.934 & 0.958 & 0.938 \\
\addlinespace
\multirow{2}{*}{CLIP ViT-B/32}   & Grayscale      & 0.938 & 0.940 & 0.225 \\
                                 & Hue rot.\ $90^\circ$ & 0.948 & 0.952 & 0.216 \\
\bottomrule
\end{tabular}
\end{table}
% note: \multirow requires \usepackage{multirow}


% ---------- Table 3: Shape vs. texture bias ----------
\begin{table}[h]
\centering
\caption{Task 1: Shape-texture cue-conflict results. }
\label{tab:t1-shape-texture}
\begin{tabular}{lccccc}
\toprule
Model & \(n_{shape}\) & \(n_{texture}\) & \(n_{other}\) & Shape bias (\%) & Coverage (\%) \\
\midrule
ResNet-50     & 81  & 74 & 85 & 52.3 & 64.6 \\
ViT-B/16       & 131 & 33 & 76 & 79.9 & 68.3 \\
CLIP ViT-B/32  & 128 & 31 & 81 & 80.5 & 66.3 \\
\bottomrule
\end{tabular}
\end{table}


% ---------- Table 4: Translation curve ----------
\begin{table}[h]
\centering
\caption{Task 1: Accuracy and prediction consistency (vs.\ zero-shift) as a function of translation magnitude.}
\label{tab:t1-translation}
\begin{tabular}{lcccccccc}
\toprule
 & \multicolumn{2}{c}{0 px} & \multicolumn{2}{c}{8 px} & \multicolumn{2}{c}{16 px} & \multicolumn{2}{c}{32 px} \\
\cmidrule(lr){2-3} \cmidrule(lr){4-5} \cmidrule(lr){6-7} \cmidrule(lr){8-9}
Model & Acc. & Cons. & Acc. & Cons. & Acc. & Cons. & Acc. & Cons. \\
\midrule
ResNet-50     & 0.972 & 1.000 & 0.965 & 0.986 & 0.967 & 0.987 & 0.959 & 0.978 \\
ViT-B/16       & 0.970 & 1.000 & 0.971 & 0.985 & 0.967 & 0.990 & 0.963 & 0.986 \\
CLIP ViT-B/32  & 0.966 & 1.000 & 0.966 & 0.973 & 0.961 & 0.963 & 0.956 & 0.966 \\
\bottomrule
\end{tabular}
\end{table}


% ---------- Table 5: Patch shuffle ----------
\begin{table}[h]
\centering
\caption{Task 1: Accuracy and consistency vs.\ clean under patch shuffling.}
\label{tab:t1-patch-shuffle}
\begin{tabular}{lcc}
\toprule
Model & Accuracy & Consistency vs.\ clean \\
\midrule
ResNet-50     & 0.888 & 0.898 \\
ViT-B/16       & 0.914 & 0.936 \\
CLIP ViT-B/32  & 0.832 & 0.818 \\
\bottomrule
\end{tabular}
\end{table}


% ---------- Table 6: Representation stability ----------
\begin{table}[h]
\centering
\caption{Task 1: Representation stability $I_T$ (cosine similarity between clean and transformed features) by intervention.}
\label{tab:t1-rep-stability}
\begin{tabular}{lccc}
\toprule
Model & Grayscale & Translation (32 px) & Patch shuffle \\
\midrule
ResNet-50     & 0.787 & 0.923 & 0.715 \\
ViT-B/16       & 0.737 & 0.951 & 0.652 \\
CLIP ViT-B/32  & 0.894 & 0.964 & 0.733 \\
\bottomrule
\end{tabular}
\end{table}

\newpage

\bibliographystyle{plainnat}  % neurips template usually sets this up already, or specifies neurips_2026
\bibliography{references}
\end{document}